% !TeX program = pdflatex
% Standalone source with an embedded bibliography.
\documentclass[11pt,a4paper,leqno]{amsart}
\usepackage[T1]{fontenc}
\usepackage{newtxtext}
\usepackage{amsmath,amsthm,amssymb,mathtools}
\usepackage{newtxmath}
\usepackage[textwidth=15.8cm,textheight=23.4cm,centering,headheight=14pt,headsep=17pt]{geometry}
\usepackage{microtype}
\usepackage{etoolbox}
\usepackage{needspace}
\usepackage[hidelinks]{hyperref}
% Counters and all mathematical environments are retained from the source.
\numberwithin{equation}{section}
\newtheoremstyle{landisplain}{7pt}{7pt}{\itshape}{}{\bfseries}{.}{.5em}{}
\newtheoremstyle{landisdefinition}{7pt}{7pt}{}{}{\bfseries}{.}{.5em}{}
\newtheoremstyle{landisremark}{6pt}{6pt}{}{}{\itshape}{.}{.5em}{}
\theoremstyle{landisplain}
\newtheorem{theorem}{Theorem}[section]
\newtheorem{lemma}[theorem]{Lemma}
\newtheorem{proposition}[theorem]{Proposition}
\newtheorem{corollary}[theorem]{Corollary}
\theoremstyle{landisdefinition}
\newtheorem{definition}[theorem]{Definition}
\theoremstyle{landisremark}
\newtheorem{remark}[theorem]{Remark}
\newcommand{\R}{\mathbb R}
\newcommand{\C}{\mathbb C}
\newcommand{\N}{\mathbb N}
\newcommand{\pa}{\partial}
\newcommand{\pb}{\bar\partial}
\newcommand{\eps}{\varepsilon}
\newcommand{\Rea}{\operatorname{Re}}
\newcommand{\Ima}{\operatorname{Im}}
\newcommand{\supp}{\operatorname{supp}}
\newcommand{\dist}{\operatorname{dist}}
\newcommand{\sgn}{\operatorname{sgn}}
\newcommand{\ord}{\operatorname{ord}}
\DeclarePairedDelimiter{\abs}{\lvert}{\rvert}
\DeclarePairedDelimiter{\norm}{\lVert}{\rVert}
\setlength{\parindent}{13pt}
\setlength{\emergencystretch}{2em}
\allowdisplaybreaks[1]
% A compact, sentence-case title and restrained journal-style hierarchy.
\makeatletter
\def\@settitle{\begin{center}%
  \normalfont\bfseries\fontsize{17}{20}\selectfont
  \@title\par
\end{center}}
\renewcommand\section{\Needspace{7\baselineskip}\@startsection{section}{1}{\z@}%
  {1.05\baselineskip \@plus .3\baselineskip \@minus .2\baselineskip}%
  {.5\baselineskip}%
  {\normalfont\large\bfseries\raggedright}}
\renewcommand\subsection{\@startsection{subsection}{2}{\z@}%
  {.75\baselineskip \@plus .2\baselineskip \@minus .1\baselineskip}%
  {.28\baselineskip}%
  {\normalfont\bfseries\raggedright}}
\renewcommand\subsubsection{\@startsection{subsubsection}{3}{\z@}%
  {.5\baselineskip \@plus .2\baselineskip}{-.5em}%
  {\normalfont\itshape}}
\makeatother
\setcounter{tocdepth}{1}
\renewcommand{\abstractname}{Abstract}
\renewcommand{\bibliofont}{\small}
\AtBeginEnvironment{thebibliography}{\setlength{\itemsep}{3pt}}
\BeforeBeginEnvironment{proposition}{\Needspace{12\baselineskip}}
\BeforeBeginEnvironment{theorem}{\Needspace{12\baselineskip}}
\BeforeBeginEnvironment{corollary}{\Needspace{12\baselineskip}}
\clubpenalty=10000
\widowpenalty=10000
\displaywidowpenalty=5000
\title[Sharp uniform biharmonic Landis estimates]
{Sharp uniform Landis estimates\\
for the biharmonic equation}
% \author{Author name}
% \address{Department, Institution, Address}
% \email{email@example.com}
\date{}
\subjclass[2020]{Primary 35B60; Secondary 35J30}
\keywords{Biharmonic equation, real-valued potential, quantitative
unique continuation, Landis estimate, vanishing order}
\hypersetup{pdftitle={Sharp uniform Landis estimates for the biharmonic equation},pdfsubject={Quantitative unique continuation; sharp uniform far-field unit-ball estimates}}

\begin{document}
\begin{abstract}
We determine the sharp uniform Landis scale for the biharmonic
equation with bounded potentials. For every sufficiently large
$R$ and every $n\ge2$, we construct smooth real-valued functions
$u_R,V_R$ satisfying $\Delta^2u_R+V_Ru_R=0$,
$\norm{V_R}_\infty\le1$ and
$u_R(0)=\norm{u_R}_\infty=1$, such that
$\norm{u_R}_{L^\infty(B_1(x_R))}
\le\exp(-cR^{4/3}\log R)$ at some point $|x_R|=R$.
Together with the lower bound of Huang, Wang and Zheng, this
establishes sharpness of the full scale $R^{4/3}\log R$ in
both real and complex scalar classes. In the plane, the potentials
are compactly supported. The construction uses an exact
biharmonic seed, adaptive removal of its poles, and a fourth-order
correction along a transverse real-part zero set. Kelvin inversion
turns the resulting annular connection into a high-order zero;
dilation produces the logarithmic factor and places the observation
ball at the prescribed radius. The examples depend on $R$:
the conclusion concerns a uniform bound over normalized pairs,
not all-direction decay of a single fixed solution.
\end{abstract}
\maketitle
\begingroup
\small
\tableofcontents
\endgroup

\section{Introduction}
\label{sec:introduction}

Quantitative unique continuation at infinity asks how small a
normalized solution of an elliptic equation can be far from the
normalization point. For bounded potentials, the answer depends on
both the operator and the scalar class. In this paper we determine
the sharp uniform scale for the biharmonic equation
\begin{equation}\label{eq:equation}
 \Delta^2u+Vu=0\quad\text{in }\R^n,\qquad n\ge2,
 \qquad \Delta=\sum_{j=1}^n\pa_j^2.
\end{equation}
For a bounded solution with $u(0)=1$, set
\begin{equation}\label{eq:mass}
 m_u(R)=\inf_{|x_0|=R}\norm{u}_{L^\infty(B_1(x_0))}.
\end{equation}
The infimum records the smallest size of the solution on a unit
ball at distance $R$ from the origin. Huang, Wang and
Zheng~\cite[Theorem~1.1]{HWZ} proved that, if
$\norm V_\infty\le1$ and $\norm u_\infty\le C_0$, then
\begin{equation}\label{eq:known-bound}
 m_u(R)\ge\exp(-CR^{4/3}\log R),\qquad R\ge R_0,
\end{equation}
where $C,R_0$ depend only on $n,C_0$. Their theorem permits
complex-valued solutions and potentials and applies to all positive
integer powers of the Laplacian. A fixed multiplicative constant
in their estimate has been absorbed into $C,R_0$.

We show that the full scale in \eqref{eq:known-bound} is optimal,
including the logarithmic factor, even when both the solution and
the potential are real. The construction attains this scale at
every prescribed sufficiently large radius. Since the resulting
pairs belong to the real class, they establish sharpness in the
larger complex classes as well. Thus real-valuedness is a strength
of the examples, not a restriction on the scope of the sharpness
conclusion.

\subsection{Main results}

\begin{theorem}\label{thm:main}
There are constants $c>0$ and $R_*>1$ such that, for every $n\ge2$
and every real number $R\ge R_*$, there exist
$u_R,V_R\in C^\infty(\R^n;\R)$ and $x_R\in\R^n$ satisfying
\begin{equation}\label{eq:normalized-examples}
 \Delta^2u_R+V_Ru_R=0,\qquad
 \norm{V_R}_\infty\le1,\qquad
 u_R(0)=\norm{u_R}_\infty=1,
\end{equation}
and
\begin{equation}\label{eq:example-smallness}
 |x_R|=R,\qquad
 \norm{u_R}_{L^\infty(B_1(x_R))}
 \le\exp(-cR^{4/3}\log R).
\end{equation}
When $n=2$, the potential $V_R$ is compactly supported, and $u_R$
is a nonzero constant outside a sufficiently large ball.
\end{theorem}

To express sharpness uniformly over the real and complex classes,
let
\[
 \mathfrak T=\{(\R,\R),(\C,\R),(\C,\C)\}.
\]
The first field specifies the values of the solution, and the
second those of the potential. For
$(\mathbb F,\mathbb K)\in\mathfrak T$ and $C_0\ge1$, let
$\mathcal C_n^{\mathbb F,\mathbb K}(C_0)$ consist of the pairs
\begin{equation}\label{eq:class}
 \begin{gathered}
 u\in L^\infty(\R^n;\mathbb F),\qquad
 V\in L^\infty(\R^n;\mathbb K),\qquad \norm V_\infty\le1,\\
 \Delta^2u+Vu=0\ \text{in distributions},\qquad
 u(0)=1,\qquad \norm u_\infty\le C_0.
 \end{gathered}
\end{equation}
The distributional formulation does not introduce an ambiguity
in the normalization. Interior elliptic regularity gives
$u\in W^{4,p}_{\mathrm{loc}}$ for every $1<p<\infty$; in
particular, for any $x\in\R^n$,
\[
 \norm u_{W^{4,p}(B_1(x))}
 \le C(n,p)\bigl(\norm u_{L^p(B_2(x))}
                 +\norm{Vu}_{L^p(B_2(x))}\bigr).
\]
Taking $p>n$ gives the continuous representative used at the
origin. The same regularity justifies smooth approximation of
compactly supported cutoff products in $W^{4,2}$ in the estimates
underlying \eqref{eq:known-bound}; no derivative of $V$ is needed.
The corresponding extremal exponent is
\begin{equation}\label{eq:extremal}
 \mathcal E_n^{\mathbb F,\mathbb K}(R;C_0)
 =\sup_{(u,V)\in\mathcal C_n^{\mathbb F,\mathbb K}(C_0)}
      [-\log m_u(R)].
\end{equation}

\begin{corollary}\label{cor:sharpness}
For every $n\ge2$ and $C_0\ge1$, there are constants $C,R_0>0$,
depending only on $n,C_0$, such that, simultaneously for all
$(\mathbb F,\mathbb K)\in\mathfrak T$,
\begin{equation}\label{eq:sharp-scale}
 cR^{4/3}\log R
 \le \mathcal E_n^{\mathbb F,\mathbb K}(R;C_0)
 \le CR^{4/3}\log R\qquad(R\ge R_0).
\end{equation}
Consequently, if $\Psi:[2,\infty)\to[0,\infty)$ satisfies
$\Psi(R)=o(R^{4/3}\log R)$, no lower bound
$m_u(R)\ge\exp(-C\Psi(R))$ can hold for all sufficiently large
$R$ uniformly over any of these classes.
\end{corollary}

\begin{proof}
The upper bound follows from \eqref{eq:known-bound}. For the
lower bound, the examples in Theorem~\ref{thm:main} lie in
$\mathcal C_n^{\R,\R}(1)$ and hence in every class under
consideration. A proposed uniform estimate with $\Psi$, applied
to the example at the same radius $R$, would give
$cR^{4/3}\log R\le C\Psi(R)$, contrary to the assumed
little-oh condition.
\end{proof}

\begin{remark}\label{rem:scope}
The order of quantifiers is
$\forall R\ge R_*\ \exists(u_R,V_R)$, not
$\exists(u,V)\ \forall R\ge R_*$. In particular, the result
does not provide a fixed pair with this smallness along an
unbounded sequence, nor a solution with this decay in every
direction. Each planar potential is compactly supported, but the
supports need not lie in a common fixed compact set. No sign
condition or uniform derivative bound is imposed on the potentials.
\end{remark}

\subsection{Background and related results}

The result sits at the intersection of three questions: local
vanishing, decay of a fixed solution at infinity, and lower bounds
uniform over a normalized family. These questions are closely
related by rescaling, but their normalizations and quantifiers
are different. We review the distinction before explaining the
construction.

\smallskip\noindent\textbf{Local unique continuation and quantitative bounds.}
For second-order equations, the qualitative theory of unique
continuation for Schr\"odinger operators was developed in particular
by Jerison and Kenig~\cite{JK}. A geometric approach through
frequency and doubling estimates was established by Garofalo and
Lin~\cite{GL86,GL87}. In the spectral setting, the work of Donnelly
and Fefferman~\cite{DF} relates the order of vanishing of an
eigenfunction to its spectral parameter. These results provide two
complementary points of view: weighted estimates propagate
nontriviality, while frequency measures the rate at which a solution
can vanish.

For equations with a variable potential, quantitative estimates must
also specify how constants depend on that potential and on the
normalization of the solution. This dependence is central to the
work of Kukavica~\cite{Kukavica}, the account of
Kenig~\cite{Kenig}, and the Carleman estimates of
Bakri~\cite{Bakri13}. It is essential to keep the normalization
visible: harmonic polynomials of arbitrarily high degree already
show that the potential norm alone cannot bound the vanishing order
of an unrestricted local solution. Bounds for regular potentials,
such as those of Bakri~\cite{Bakri12} and Zhu~\cite{Zhu16},
measure additional coefficient regularity. They therefore concern
a different uniformity question from one in which only
$\norm V_\infty$ is controlled.

\smallskip\noindent\textbf{From Landis' question to a unit-ball estimate.}
The classical Landis question concerns whether a nonzero solution
of a Schr\"odinger equation with bounded potential can decay faster
than an exponential at infinity. In the complex-valued class,
Meshkov~\cite{Meshkov} constructed a bounded potential and a
nontrivial solution in the plane with decay of order
$\exp(-c|x|^{4/3})$. His work identifies the power $4/3$ as the
critical one for that complex-valued decay problem. Motivated by
Anderson localization, Bourgain and Kenig~\cite{BK} obtained the
quantitative unit-ball lower bound with exponent
$CR^{4/3}\log R$ for bounded normalized solutions of the
second-order equation. Extensions of unique continuation at
infinity to other coefficient regimes include
Cruz-Sampedro~\cite{CS} for complex potentials, Davey~\cite{Davey}
for magnetic Schr\"odinger operators, and Lin and Wang~\cite{LW}
for general second-order elliptic equations.

The logarithm distinguishes the unit-ball problem from the
critical-power question. Under a large-scale rescaling, a fixed
observation ball becomes a small ball; a power-law lower bound
on that ball then produces a logarithmic factor. Conversely, an
example bounded above by $\exp(-cR^{4/3})$ does not exclude a
uniform lower bound $\exp(-CR^{4/3})$ with a larger constant
$C$. Establishing that the logarithm is necessary requires a
smaller example, with the global normalization retained.

\smallskip\noindent\textbf{The real-valued problem and the role of dimension.}
In the plane, real-valuedness leads to a different second-order
theory. Kenig, Silvestre and Wang~\cite{KSW} proved a quantitative
Landis estimate under a nonnegative Schr\"odinger potential; the
sign convention here is $-\Delta w+Ww=0$ with $W\ge0$.
Davey, Kenig and Wang~\cite{DKW} extended this line of analysis
to variable-coefficient planar equations. The sign-free real
problem was subsequently treated by Logunov, Malinnikova,
Nadirashvili and Nazarov~\cite{LMNN}. For the comparison relevant
here, their Theorem~1.2 implies that a real entire solution of
$\Delta w+Ww=0$ in $\R^2$, with $W$ real,
$\norm W_\infty\le1$ and $w(0)=\norm w_\infty=1$, satisfies
\[
 \inf_{|x_0|=R}\norm{w}_{L^\infty(B_1(x_0))}
 \ge \exp\bigl(-CR(\log R)^{3/2}\bigr)
\]
for all sufficiently large $R$. This follows by applying their
local theorem on $B_{4R}$ with its radius parameter equal to
$2R$. Their theorem on decay in every direction is a separate
statement; it should not be identified with this unit-ball estimate.

The dimension restriction is substantive. In a recent preprint,
Frank and Ivanisvili~\cite[Theorem~1]{FI} construct, in every
dimension $n\ge3$, a nonzero smooth real solution with a smooth
bounded real potential and all-direction decay
$\exp(-c|x|^{4/3})$. Thus a general contrast between real
second-order and fourth-order equations would be inaccurate.
The contrast established here is already present in the plane:
real-valuedness improves the second-order unit-ball bound beyond
the power $4/3$, whereas for the biharmonic operator the full
scale $R^{4/3}\log R$ remains necessary. Our construction then
extends to all dimensions $n\ge2$.

A related, but distinct, local result is due to Feng, Gan, Jiang
and Lin~\cite[Theorem~1.1]{FGJL}. For every sufficiently large
integer $k$, they construct real Dirichlet solutions in
$B_3\subset\R^3$ with vanishing order $k$ and potential norm
at most $Ck^{3/2}$. This is a fixed-domain second-order
construction. It does not directly supply the bounded entire
solutions, the maximum normalization, or the prescribed distant
observation ball required here. Neither this result nor the
second-order examples above are used as inputs to our proof.

\smallskip\noindent\textbf{Powers of the Laplacian.}
For higher-order equations, unique continuation must be studied
with the structure of the principal part in view. Strong unique
continuation for powers of the Laplacian was investigated by
Colombini and Grammatico~\cite{CG} and Lin~\cite{Lin07}.
Lin, Nagayasu and Wang~\cite{LNW} developed quantitative
estimates with singular lower-order coefficients, including
three-sphere and doubling inequalities. Further estimates for
vanishing orders, with explicit coefficient dependence, were
obtained by Zhu~\cite{Zhu16,Zhu18}. More recently, Liu, Tian
and Yang~\cite{LTY} studied quantitative unique continuation
for a bi-Laplacian equation with a $W^{1,\infty}$ potential.
These derivative-sensitive estimates do not give a uniform
improvement over a class in which no derivative of the potential
is bounded.

On the counterexample side, Fu, Zhang and Zuazua~\cite{FZZ}
constructed complex-valued biharmonic examples with decay of
order $\exp(-c|x|^{4/3})$ in connection with the optimality
of observability inequalities for plate equations. On the
lower-bound side, Huang, Wang and Zheng~\cite{HWZ} established
\eqref{eq:known-bound}. The present construction matches that
lower bound in the smallest scalar class, with smooth real
solutions and smooth real potentials. It addresses the logarithmic
loss that is not settled by a critical-power decay example.
Smoothness here is qualitative: the potentials have a uniform
$L^\infty$ bound, but no uniform derivative bound or common
compact support is asserted.

\subsection{Outline of the construction}

The argument begins in the plane, where the identity
$\Delta^2=16\pa^2\pb^2$ provides exact biharmonic products.
At degrees $k=q^2$, a polynomial with separated roots determines
a complex seed whose inner and outer models are
$A_k|z|^2$ and $|z|^2z^{-k}$. Matching on fixed collars
introduces only an exponentially small residual. The poles at
the roots must then be removed before a real potential can be
defined.

Taking a real part is not sufficient: the numerator
$\Delta^2\Rea F$ need not vanish on the zero set of
$\Rea F$. Section~\ref{sec:repair} resolves this obstruction
where that zero set is transverse. With
$h=\Rea F/|F|$, a correction of order $h^4$ preserves the
first three jets and cancels the fourth-order residual on $h=0$.
A quantitative gradient bound then gives smooth division and a
bounded real quotient. In the root cores, by contrast, the
replacement remains exactly biharmonic. Possible critical zeros
there require no division. The argument uses both mechanisms,
not the Taylor correction alone.

Sections~\ref{sec:seed} and~\ref{sec:adaptive} arrange the
geometry needed by this correction. The two replacement radii
$\eps_0k^{-1/2}$ and $\eps_1k^{-2/3}$ distinguish roots
where the holomorphic phase dominates from those where the
replacement phase must dominate. This avoids cancellation of
phase gradients and yields the decisive bound $D_\xi\le Ck^3$.
Disjoint residual supports then give a real annular connection
with $\norm{Q_k}_\infty\le Ck^3$, without a loss proportional
to the number of roots
(Proposition~\ref{prop:annular}).

Finally, Section~\ref{sec:kelvin} applies Kelvin inversion. The
outer model becomes a harmonic polynomial of degree $k$ near
zero, while the inner model becomes a nonzero constant near
infinity. A dilation of order $k^{3/4}$ normalizes the potential.
On the observation ball, the polynomial contributes the factor
$\lambda^{-k}=\exp(-k\log\lambda)$. Choosing
$k\asymp R^{4/3}$ and $\lambda\asymp R$ gives precisely
the scale $R^{4/3}\log R$. Translation from a point of maximum
modulus and the final choice of dilation place the center at
exactly the prescribed radius, not merely along a sequence.

\section{A real-part correction}
\label{sec:repair}

The correction preserves the real part through third order on
its zero set, while its fourth normal jet cancels the biharmonic
residual there through the identity
$\Delta^2(h^4a)|_{h=0}=24a|\nabla h|^4$. When phase transversality
makes $\nabla h$ non-degenerate, this cancellation permits smooth
division by the corrected real function, without requiring that
function to be nonvanishing. The estimate below quantifies this
step; exact biharmonicity will handle the critical zeros outside
the correction regions.

\subsection{Notation and biharmonic products}

We identify $\R^2$ with $\C$ and write $z=x_1+ix_2$.
The notation $|D^jF|$ denotes the maximum of $|\pa^\alpha F|$
over real multiindices $\alpha$ with $|\alpha|=j$; gradient norms
are Euclidean. Unless indicated otherwise, constants depend only
on fixed radii, cutoff functions and the stated finite derivative
orders, never on $q,k$ or an individual root. We write $X\asymp Y$
when $cY\le X\le CY$ with such fixed positive constants.

Set
\[
 \pa=\frac12(\pa_{x_1}-i\pa_{x_2}),\qquad
 \pb=\frac12(\pa_{x_1}+i\pa_{x_2}),\qquad
 \Delta^2=16\pa^2\pb^2.
\]
If $H$ is holomorphic and $G$ is smooth, then
\begin{equation}\label{eq:prod}
 \Delta^2(HG)=16\left(H''\pb^2G+2H'\pa\pb^2G
                                  +H\pa^2\pb^2G\right).
\end{equation}
Since $\pb H=0$, the identity follows by applying $\pb^2$ to
$G$ and then using the product rule for $\pa^2$. In particular,
$\pb^2G=0$ implies $\Delta^2(HG)=0$. This exact cancellation
will preserve biharmonicity in the replacement cores.

\subsection{Correction along a transverse zero set}

For a nonvanishing complex-valued smooth function $F$, set
\[
 \nabla\arg F:=\Ima\left(\frac{\nabla F}{F}\right).
\]
This definition does not require a single-valued choice of
argument on the domain.

\begin{lemma}[Real-part correction]\label{lem:repair}
Let $\Omega\subset\R^2$ be open and let
$F\in C^\infty(\Omega;\C)$ be nonvanishing. Put
\[
 A=|F|,\qquad f=\Rea F,\qquad h=f/A,\qquad E=\Delta^2f.
\]
Let $L\ge1$ and $D>0$. Assume that, for fixed positive constants
$b,c,C_*$ and $\{C_j\}_{j=0}^9$, the relative derivative bounds
and phase estimate
\begin{align}
 |D^jF|&\le C_jL^jA &&(0\le j\le9),\label{eq:Fder}\\
 |D^jE|&\le C_jDL^jA &&(0\le j\le5),\label{eq:Eder}\\
 |\nabla\arg F|&\ge cL
 &&\text{on }\{x\in\Omega:|h(x)|\le1/2\},\label{eq:phase}
\end{align}
hold, and that the residual has the size and support margin
\begin{equation}\label{eq:margin}
 D\le C_*L^4,\qquad
 \supp E\Subset\Omega,\qquad
 \dist(\supp E,\R^2\setminus\Omega)\ge b/L.
\end{equation}
Then there exist $U\in C^\infty(\Omega;\R)$ and
$Q\in C_c^\infty(\Omega;\R)$ such that
\begin{equation}\label{eq:repair-result}
 \Delta^2U+QU=0,\qquad \norm Q_\infty\le CD,
 \qquad U=f\ \text{off }\supp E.
\end{equation}
The ratio $U/f$ extends smoothly with values in $[1/2,3/2]$,
and $U-f$ vanishes to fourth order on $\{f=0\}$. Here $C$
depends only on $b,c,C_*$ and the constants in
\eqref{eq:Fder}--\eqref{eq:Eder}.
\end{lemma}

\begin{proof}
We first record the relative derivative bounds
\begin{equation}\label{eq:Ah-bounds}
 |D^jA|\le C_j'L^jA,\qquad
 |D^j(A^{-1})|\le C_j'L^jA^{-1},\qquad
 |D^jh|\le C_j'L^j\quad(0\le j\le9).
\end{equation}
These follow inductively from $A^2=F\bar F$: the highest
multiindex derivative of $A$ has coefficient $2A$, and each
lower-order product has the required homogeneity. Differentiating
$AA^{-1}=1$ gives the estimate for $A^{-1}$, after which the
product rule applied to $h=fA^{-1}$ gives the estimate for $h$.
The constants do not require a positive lower bound for $A$.
Writing $F=Ae^{i\phi}$ locally, we have $h=\cos\phi$;
therefore \eqref{eq:phase} yields
\begin{equation}\label{eq:h-gradient}
 |\nabla h|=\sqrt{1-h^2}\,|\nabla\phi|
 \ge c_1L\quad\text{on }|h|\le1/2,
 \qquad c_1=\sqrt3c/2.
\end{equation}

\smallskip\noindent\emph{Choice of the correction.}
Fix $\chi\in C_c^\infty((-2,2);\R)$ with $0\le\chi\le1$ and
$\chi=1$ on $[-1,1]$. Write $|E|\le C_EDA$ and set
$K=C_EC_*/(3c_1^4)$. Choose once and for all
\begin{equation}\label{eq:delta-choice}
 0<\delta\le\min\left\{\frac18,\frac{c_1b}{4},(2K)^{-1/3}\right\}.
\end{equation}
Define
\begin{equation}\label{eq:correction}
 B=\chi(h/\delta)\frac{h^4E}{24|\nabla h|^4},\qquad U=f-B.
\end{equation}
The quotient is evaluated only on $|h|<2\delta$, and $B$ is
set to zero elsewhere. The gradient bound
\eqref{eq:h-gradient} and the compact support of $\chi$ make
this definition smooth, with $\supp B\subset\supp E$.
For $h\ne0$, the relative correction satisfies
\[
 \left|\frac Bf\right|
 \le\frac{C_ED(2\delta)^3}{24c_1^4L^4}
 \le K\delta^3\le\frac12.
\]
The same ratio extends smoothly across $h=0$. Consequently,
\begin{equation}\label{eq:comparison}
 U=Ah\Theta,\qquad
 \Theta=1-\chi(h/\delta)\frac{h^3E}{24A|\nabla h|^4},
 \qquad \frac12\le\Theta\le\frac32,
\end{equation}
with the quotient in $\Theta$ interpreted by the same cutoff rule.
Thus $U$ and $f$ have exactly the same zero set, and
\begin{equation}\label{eq:transverse-U}
 \nabla U=A\nabla h\ne0\quad\text{on }\{h=0\}.
\end{equation}

\smallskip\noindent\emph{Cancellation on the zero set.}
For any smooth functions $p,a$,
\begin{align*}
 \Delta^2(pa)={}&a\Delta^2p+p\Delta^2a
       +2\Delta p\,\Delta a
       +4\nabla a\cdot\nabla\Delta p\\
       &+4\nabla p\cdot\nabla\Delta a
       +4\sum_{i,j=1}^2p_{ij}a_{ij}.
\end{align*}
For $p=h^4$, all derivatives of $p$ of order at most three vanish
on $h=0$. Moreover,
\[
 \Delta(h^4)=12h^2|\nabla h|^2+4h^3\Delta h,
 \qquad
 \left.\Delta^2(h^4)\right|_{h=0}=24|\nabla h|^4.
\]
Thus every term involving a derivative of $a$ vanishes on the
zero set, leaving precisely
\begin{equation}\label{eq:jet}
 \left.\Delta^2(h^4a)\right|_{h=0}=24a|\nabla h|^4.
\end{equation}
Since $\chi=1$ near $h=0$, applying this identity to
\eqref{eq:correction} cancels the residual there:
\begin{equation}\label{eq:N-zero}
 N:=\Delta^2U=0\quad\text{on }\{h=0\}.
\end{equation}

\smallskip\noindent\emph{Derivative estimates.}
On $|h|<2\delta$, put $v=L^{-1}\nabla h$. By
\eqref{eq:Ah-bounds}--\eqref{eq:h-gradient},
\[
 |v|\ge c_1,\qquad |D^rv|\le C_rL^r\quad(0\le r\le5).
\]
Differentiating $|\nabla h|^{-4}=L^{-4}|v|^{-4}$ therefore yields
\begin{equation}\label{eq:inverse-gradient}
 |D^r(|\nabla h|^{-4})|\le C_rL^{r-4}
 \qquad(0\le r\le5).
\end{equation}
This step uses derivatives of $h$ only through order six.
Likewise,
\[
 |D^r[\chi(h/\delta)h^4]|\le C_{r,\delta}L^r.
\]
In a product term with $r+s+t=j$, the three factors in
\eqref{eq:correction} are bounded by $CL^r$, $CDL^sA$ and
$CL^{t-4}$. The product rule thus gives
\begin{equation}\label{eq:Bder}
 |D^jB|\le CD L^{j-4}A\qquad(0\le j\le5).
\end{equation}
Since $\delta$ is fixed, its inverse powers enter only the
constants. Consequently,
\begin{equation}\label{eq:Nd}
 |N|\le CDA,\qquad |\nabla N|\le CDLA,
 \qquad \supp N\subset\supp E.
\end{equation}
The last support assertion uses the locality of differentiation.

\smallskip\noindent\emph{Division estimate near the zero set.}
To divide by $U$, we must recover the factor $|h|$ missing
from the first estimate in \eqref{eq:Nd}. For
$x\in\supp E$ with $0<|h(x)|<\delta$, set
$\sigma=\sgn(h(x))$ and follow the normalized gradient flow
\begin{equation}\label{eq:flow}
 \dot\gamma(t)=-\sigma\frac{\nabla h(\gamma(t))}
                                  {|\nabla h(\gamma(t))|^2},
 \qquad \gamma(0)=x.
\end{equation}
As long as the solution remains in $|h|\le\delta$,
\[
 h(\gamma(t))=h(x)-\sigma t,
 \qquad |\dot\gamma(t)|\le(c_1L)^{-1}.
\]
Up to time $T:=|h(x)|$, the path has length at most
$T/(c_1L)\le b/(4L)$. The support margin
$\dist(x,\R^2\setminus\Omega)\ge b/L$ therefore keeps it in
a compact subset of $\Omega$, where the gradient is bounded away
from zero along the flow. The local ODE solution extends to time
$T$ and reaches the zero set: $h(\gamma(T))=0$.

The relative amplitude bound in \eqref{eq:Ah-bounds} gives
\[
 \left|\frac{d}{dt}\log A(\gamma(t))\right|
 \le \frac{C_A}{c_1},\qquad
 e^{-C_A\delta/c_1}A(x)\le A(\gamma(t))
                         \le e^{C_A\delta/c_1}A(x).
\]
Using \eqref{eq:N-zero} at the endpoint and integrating
\eqref{eq:Nd}, we obtain
\begin{equation}\label{eq:N-division}
 |N(x)|\le\int_0^T|\nabla N(\gamma(t))|\,|\dot\gamma(t)|\,dt
 \le CDA(x)|h(x)|.
\end{equation}
At $h(x)=0$, the same estimate follows from the cancellation;
off $\supp E$, it follows from $N=0$. On $|h|\ge\delta$,
the comparison \eqref{eq:comparison} and the residual estimate
\eqref{eq:Nd} give the quotient bound without a flow argument.
Combining the two regions, we obtain
\begin{equation}\label{eq:N-U-bound}
 |N|\le CD|U|\quad\text{on }\Omega.
\end{equation}

\smallskip\noindent\emph{Smooth extension of the quotient.}
Define $Q=-N/U$ where $U\ne0$. At a zero, use smooth local
coordinates $(s,t)$ with $s=h$, which exist by
\eqref{eq:h-gradient}. In these coordinates,
\[
 N(s,t)=s\int_0^1\pa_sN(\tau s,t)\,d\tau=:s\widetilde N(s,t),
 \qquad U(s,t)=sA(s,t)\Theta(s,t).
\]
Hence $Q=-\widetilde N/(A\Theta)$ is smooth. The local
extensions agree on overlaps because they agree away from $h=0$,
a dense subset of each chart. By continuity,
\eqref{eq:N-U-bound} gives $\norm Q_\infty\le CD$ also on
the zero set. Moreover, $Q$ vanishes off $\supp E$ and is
therefore compactly supported in $\Omega$. The construction is
real-valued throughout, and \eqref{eq:correction} gives the
claimed fourth-order vanishing.
\end{proof}

\begin{remark}\label{rem:repair-scope}
Smooth division requires transversality of $U$, or equivalently
of $h$, not of $N=\Delta^2U$; the numerator need only vanish
on the hypersurface $h=0$. For the quantitative estimate, the phase
bound must hold throughout a neighborhood of the residual support
with the margin in \eqref{eq:margin}. A bound restricted to points
where $E\ne0$ would not justify the flow argument. No such
condition is needed at critical zeros lying in open regions of
exact biharmonicity.
\end{remark}

\section{The polynomial seed and exact matching}
\label{sec:seed}

We next construct the complex seed to which the local correction
will be applied. Its roots must be sufficiently separated for
independent replacements, while its inner and outer collars must
match the exact models with negligible residual. We use degrees
$k=q^2$ with $q\in\N$; the choice in
Section~\ref{sec:kelvin} will show that these degrees suffice
for every prescribed large radius. Set
\begin{equation}\label{eq:P}
 \rho_j=2+\frac{j}{q+1}\quad(1\le j\le q),\qquad
 P(z)=\prod_{j=1}^q(z^q-\rho_j^q).
\end{equation}
Then $P$ is monic of degree $k$, $P(0)\in\R\setminus\{0\}$, and
its root set $\Lambda$ consists of $k$ simple points in $2<|z|<3$.
For roots on different circles, the radial separation is at least
$(q+1)^{-1}$; on one circle, the minimum distance is
$2\rho_j\sin(\pi/q)$. Thus, for $q\ge2$,
\begin{equation}\label{eq:root-separation}
 |\xi-\eta|\ge c_*k^{-1/2}\qquad(\xi\ne\eta,\ \xi,\eta\in\Lambda)
\end{equation}
with a fixed $c_*>0$.

Choose $\eps_0>0$ so small that the closed disks of radius
$3\eps_0k^{-1/2}$ centered at the roots are pairwise disjoint and lie
in $7/4<|z|<13/4$. For each $\xi\in\Lambda$, put
\begin{equation}\label{eq:H}
 H_\xi(z)=\frac{z}{P(z)/(z-\xi)},\qquad
 \ell_\xi(z)=\frac{H_\xi'(z)}{H_\xi(z)}.
\end{equation}
The factor $H_\xi$ is holomorphic and nonvanishing on its disk;
it isolates the pole at $\xi$ from the contribution of the
remaining roots.

\subsection{Logarithmic derivatives}

\begin{lemma}\label{lem:log}
Uniformly in $\xi$ and on $B_{3\eps_0k^{-1/2}}(\xi)$,
\begin{align}
 |\ell_\xi|&\le Ck,\label{eq:ell1}\\
 |\ell_\xi'|&\le Ck\log k,\label{eq:ell2}\\
 |\pa_z^{j-1}\ell_\xi|&\le C_jk^{j/2}
           \qquad(3\le j\le9).\label{eq:ellj}
\end{align}
\end{lemma}

\begin{proof}
Root separation first gives a uniform packing estimate: the disks
of radius a fixed multiple of $k^{-1/2}$ centered at the roots
are disjoint, so
\begin{equation}\label{eq:packing}
 \#(\Lambda\cap B_t(z))\le C(1+kt^2).
\end{equation}
For $z$ in the disk associated with $\xi$, all other roots are a
distance at least $d_0=ck^{-1/2}$ away, after decreasing $\eps_0$
if needed. All roots and these disks lie in a fixed bounded set.
Decomposing that set into rings
$2^\nu d_0\le|z-\eta|<2^{\nu+1}d_0$ gives
\[
 \sum_{\eta\ne\xi}|z-\eta|^{-s}
 \le Ck\sum_{\nu=0}^{O(\log k)}(2^\nu d_0)^{2-s}
 \le
 \begin{cases}
 Ck,&s=1,\\
 Ck\log k,&s=2,\\
 C_sk^{s/2},&s>2.
 \end{cases}
\]
The logarithmic loss is confined to the borderline case $s=2$.
Now use
\[
 \ell_\xi(z)=\frac1z-\sum_{\eta\ne\xi}\frac1{z-\eta}
\]
and differentiate. Since $z$ stays a fixed positive distance
from zero, the root sums give all the required bounds; explicitly,
\[
 \pa_z^{j-1}\ell_\xi
 =(-1)^{j-1}(j-1)!
   \left(z^{-j}-\sum_{\eta\ne\xi}(z-\eta)^{-j}\right).
\]
\end{proof}

\subsection{Exponential matching to the exact models}

The same arrangement of roots also makes the seed exponentially
close to its exact models on fixed collars. Define
\[
 E_-(z)=P(z)/P(0)-1,\qquad E_+(z)=P(z)/z^k-1,
\]
with $E_+$ considered away from zero. We use the fixed enlarged
collars
\begin{equation}\label{eq:collar-domains}
 \Omega_-:=\{9/8<|z|<13/8\},\qquad
 \Omega_+:=\{27/8<|z|<31/8\}.
\end{equation}
There are fixed $0<\vartheta<1$ and constants $C_j$ such that
\begin{equation}\label{eq:collarexp}
 \begin{aligned}
 |D^jE_-|&\le C_jq\vartheta^q&&\text{on }|z|\le13/8,\\
 |D^jE_+|&\le C_jq\vartheta^q&&\text{on }|z|\ge27/8,
 \end{aligned}
 \qquad 0\le j\le9.
\end{equation}
Indeed,
\begin{equation}\label{eq:collar-products}
 \frac{P(z)}{P(0)}=\prod_{j=1}^q\left(1-(z/\rho_j)^q\right),\qquad
 \frac{P(z)}{z^k}=\prod_{j=1}^q\left(1-(\rho_j/z)^q\right).
\end{equation}
On the larger regions $|z|\le7/4$ and $|z|\ge13/4$, respectively,
every power has modulus at most $\vartheta_0^q$ for a fixed
$\vartheta_0<1$. The product error is bounded by
$\exp(q\vartheta_0^q)-1\le Cq\vartheta_0^q$ for large $q$.
Cauchy estimates on disks of a fixed radius, contained in these larger
regions, give \eqref{eq:collarexp}. The same bounds hold for real
partial derivatives because the errors are holomorphic.

Choose fixed smooth radial cutoffs taking values in $[0,1]$ with
\begin{align*}
 &\eta_-(r)=1\quad(r\le5/4),\qquad
 \eta_-(r)=0\quad(r\ge3/2),\\
 &\eta_+(r)=0\quad(r\le7/2),\qquad
 \eta_+(r)=1\quad(r\ge15/4).
\end{align*}
Set
\begin{equation}\label{eq:seed}
 S=1+\eta_-E_-+\eta_+E_+,\qquad
 F^0(z)=\frac{|z|^2}{P(z)}S(z)\quad(z\notin\Lambda).
\end{equation}
The outer cutoff term is set to zero near the origin. By
\eqref{eq:collarexp}, we have $|S-1|\le1/2$ everywhere for
large $k$, whereas $S=1$ on every root disk. Thus the collar
modifications leave the local pole structure unchanged and give
the exact models
\begin{equation}\label{eq:models}
 \begin{aligned}
 F^0(z)&=P(0)^{-1}|z|^2&&(|z|\le5/4),\\
 F^0(z)&=|z|^2z^{-k}&&(|z|\ge15/4).
 \end{aligned}
\end{equation}

Away from the roots, the only biharmonic residual is confined
to the two closed transition collars
\[
 \mathcal A_-:=\{5/4\le|z|\le3/2\},\qquad
 \mathcal A_+:=\{7/2\le|z|\le15/4\}.
\]
Their distances from the boundaries of $\Omega_-$ and
$\Omega_+$, respectively, are at least $1/8$. This fixed
margin will permit the local correction without changing either
exact model.

\subsection{Relative bounds and phase on the collars}

On each enlarged collar, $F^0$ is nonzero and satisfies
\begin{align}
 |D^jF^0|&\le C_jk^j|F^0| &&(0\le j\le9),\label{eq:collarF}\\
 |D^j\Delta^2F^0|&\le C_jk^{j+2}q\vartheta^q|F^0|
                                     &&(0\le j\le5).\label{eq:collarres}
\end{align}
To verify the second estimate, write $H=z/P$ and $G=\bar zS$ in
\eqref{eq:prod}. Since the collars stay a fixed distance from all roots
and zero, logarithmic differentiation gives
$|H^{(j)}|\le C_jk^j|H|$ for the required orders. Moreover,
\[
 \pb^2(\bar zS)=2\pb S+\bar z\pb^2S.
\]
Every nonzero term, including those obtained by further
differentiation, contains a derivative of the exponentially small
function $S-1$. Applying \eqref{eq:prod} and then at most
five additional derivatives gives \eqref{eq:collarres}.
At most nine derivatives fall on the collar error, while the
initial product identity places at most two derivatives on $H$.
The relative cost is therefore $k^{j+2}$. Finally, $|G|$ is
bounded above and below by fixed positive constants on the collars,
so the estimate may be divided by $|F^0|=|H||G|$.

On $\Omega_-$, the formula can be written as
\[
 F^0=P(0)^{-1}|z|^2(1+T_-),\qquad
 T_-=(\eta_--1)\frac{E_-}{1+E_-}.
\]
Since $P(0)^{-1}$ is real and $|T_-|\le Cq\vartheta^q$,
\begin{equation}\label{eq:inner-real}
 |\Rea F^0|\ge\frac12|F^0|\quad\text{on }\Omega_-
\end{equation}
for large $k$. The real part is bounded away from zero relative
to its amplitude, so no zero-set correction is needed on the
inner collar. On $\Omega_+$,
\[
 F^0=|z|^2z^{-k}(1+T_+),\qquad
 T_+=(\eta_+-1)\frac{E_+}{1+E_+},\qquad
 |D^jT_+|\le C_jq\vartheta^q\quad(j\le9).
\]
Therefore
\begin{equation}\label{eq:outerphase}
 |\nabla\arg F^0|
 \ge\frac{k}{|z|}-\left|\frac{\nabla T_+}{1+T_+}\right|
 \ge ck\quad\text{on }\Omega_+.
\end{equation}
Thus the inner collar permits direct division, while the outer
collar has the phase transversality needed for correction. Both
collars in \eqref{eq:collar-domains} are disjoint from the root
disks, and all preceding constants are independent of the root
replacements. It remains to remove the poles without losing this
separation or exceeding the potential budget.

\section{Adaptive removal of the singularities}
\label{sec:adaptive}

The replacement must do more than remove a pole: it must keep
the real-part zero set transverse wherever a residual is created.
We choose the radius so that either the holomorphic phase or the
local replacement phase dominates the other. This two-scale choice
prevents phase cancellation and keeps the residual within the
required budget.

Fix a smooth function $\chi_0:[0,\infty)\to[0,1]$ such that
$\chi_0=0$ on $[0,1]$ and $\chi_0=1$ on $[2,\infty)$.
For $\xi\in\Lambda$, write $w=z-\xi$ and define
\begin{equation}\label{eq:G}
 G_{\xi,\eps}(z)
 =\chi_0(|w|/\eps)\frac{\bar z}{w}
 +(1-\chi_0(|w|/\eps))\frac{\bar\xi}{\eps^2}\bar w.
\end{equation}
The first summand is defined to be zero near $w=0$.

\begin{lemma}\label{lem:G}
On $\Omega_{\xi,\eps}:=\{\eps/2<|z-\xi|<3\eps\}$, uniformly for
$2\le|\xi|\le3$ and all sufficiently small $\eps>0$,
\begin{align}
 c\eps^{-1}\le |G_{\xi,\eps}|&\le C\eps^{-1},\label{eq:Gsize}\\
 |D^jG_{\xi,\eps}|&\le C_j\eps^{-j}|G_{\xi,\eps}|
                      \qquad(0\le j\le9),\label{eq:Gder}\\
 \nabla\arg G_{\xi,\eps}&=-\nabla\arg w+O(1).\label{eq:Gphase}
\end{align}
The $O(1)$ is a real vector with a bound independent of $\xi,\eps$.
Whenever $H_\xi$ is defined, the product $H_\xi G_{\xi,\eps}$ is
biharmonic on $|w|<\eps$ and on $|w|>2\eps$.
\end{lemma}

\begin{proof}
Write $w=se^{i\theta}$ and abbreviate $\chi_0(s/\eps)$ by $\chi_0$.
Then
\[
 G=\bar\xi e^{-i\theta}B(s)+\chi_0e^{-2i\theta},\qquad
 B(s)=\frac{\chi_0}{s}+(1-\chi_0)\frac{s}{\eps^2}.
\]
On $\eps/2<s<3\eps$, the positive quantity $B$ lies between
$c/\eps$ and $C/\eps$. Hence
\begin{equation}\label{eq:G-factor}
 G=\bar\xi e^{-i\theta}B(s)(1+T),\qquad
 T=\frac{\chi_0e^{-i\theta}}{\bar\xi B(s)},
 \qquad |T|\le C\eps,\quad |\nabla T|\le C.
\end{equation}
To estimate the derivative in \eqref{eq:G-factor}, write
$B(s)=\eps^{-1}b(s/\eps)$. The positive function $b$ and its
reciprocal have uniformly bounded derivatives on $[1/2,3]$.
Thus $T$ is $\eps$ times a smooth profile of $w/\eps$, which
gives $|\nabla T|\le C$. For small $\eps$, the factorization
\eqref{eq:G-factor} yields \eqref{eq:Gsize}. Since the radial
factor $B$ is real and positive, it contributes no phase:
\[
 \nabla\arg G=-\nabla\theta+
                  \Ima\frac{\nabla T}{1+T},
\]
which proves \eqref{eq:Gphase}.

For the higher derivative estimates, put $w=\eps\zeta$ and write
\[
 G=\frac{\bar\xi}{\eps}
       \left(\frac{\chi_0(|\zeta|)}{\zeta}
                         +(1-\chi_0(|\zeta|))\bar\zeta\right)
       +\chi_0(|\zeta|)\frac{\bar\zeta}{\zeta}.
\]
On the fixed ring $1/2<|\zeta|<3$, all derivatives through order nine
of the two bracketed profiles are bounded. The chain rule gives
$|D^jG|\le C_j\eps^{-j-1}$, which, together with the lower bound in
\eqref{eq:Gsize}, yields \eqref{eq:Gder}.
In the core and outside the transition ring, the exact profiles
are $(\bar\xi/\eps^2)\bar w$ and $\bar z/w$, respectively.
Both are annihilated by $\pb^2$ on their domains, so
\eqref{eq:prod} proves exact biharmonicity in both regions.
\end{proof}

The radius is now chosen from the strength of the holomorphic
phase. Define $a_\xi=|\ell_\xi(\xi)|$. With $\eps_0$ fixed,
choose a small constant $\eps_1>0$ and set
\begin{equation}\label{eq:adaptive}
 \eps_\xi=
 \begin{cases}
 \eps_0k^{-1/2},&a_\xi\ge k^{2/3},\\
 \eps_1k^{-2/3},&a_\xi<k^{2/3}.
 \end{cases}
\end{equation}
For all sufficiently large $k$, every disk $B_{3\eps_\xi}(\xi)$ is
contained in its original separation disk. In these disks replace
$F^0$ by $H_\xi G_{\xi,\eps_\xi}$, and retain $F^0$ elsewhere.
Since $S=1$ on the root disks, the formulas agree on the open
rings $2\eps_\xi<|z-\xi|<3\eps_\xi$. They therefore define
a global smooth complex-valued function $F$, with no residual
introduced at the gluing boundaries.

\begin{proposition}\label{prop:patch}
For each enlarged root ring
$\Omega_\xi:=\{\eps_\xi/2<|z-\xi|<3\eps_\xi\}$, the
function $F$ is nonvanishing and there are $L_\xi\ge1$, $D_\xi>0$
such that
\begin{align}
 |D^jF|&\le C_jL_\xi^j|F| &&(0\le j\le9),\label{eq:patchF}\\
 |D^j\Delta^2F|&\le C_jD_\xi L_\xi^j|F|
                                      &&(0\le j\le5),\label{eq:patchE}\\
 |\nabla\arg F|&\ge cL_\xi,
 \qquad D_\xi\le CL_\xi^4,
 \qquad D_\xi\le Ck^3.\label{eq:patchphase}
\end{align}
The constants are uniform in $k$ and $\xi$. The residual is supported
in $\{\eps_\xi\le|z-\xi|\le2\eps_\xi\}$, at distance at least
$b/L_\xi$ from $\pa\Omega_\xi$, for a fixed $b>0$.
\end{proposition}

\begin{proof}
We first separate the phase estimate from the derivative costs.
For a nonvanishing holomorphic function $H$, with $\ell=H'/H$,
\begin{equation}\label{eq:holomorphic-phase}
 \nabla\arg H=(\Ima\ell,\Rea\ell),\qquad
 |\nabla\arg H|=|\ell|.
\end{equation}
Moreover, $H^{(j)}/H$ is a polynomial in
$\ell,\ell',\ldots,\ell^{(j-1)}$, homogeneous of weight $j$
when $\ell^{(r)}$ is assigned weight $r+1$. Consequently, the
bounds $|\ell^{(r)}|\le C_rL^{r+1}$ for $0\le r\le8$
imply
\begin{equation}\label{eq:H-relative}
 |H^{(j)}|\le C_jL^j|H|\qquad(0\le j\le9).
\end{equation}

\smallskip\noindent\emph{Case A: $a_\xi\ge k^{2/3}$.}
Set $L_\xi=a_\xi$ and $\eps_\xi=\eps_0k^{-1/2}$.
Lemma~\ref{lem:log} gives, throughout the corresponding disk,
\[
 |\ell_\xi(z)-\ell_\xi(\xi)|
 \le Ck\log k\,\eps_\xi
 \le C\eps_0k^{1/2}\log k=o(k^{2/3}).
\]
For all sufficiently large $k$, with a threshold independent
of $\xi$, we therefore have
\[
 \frac34L_\xi\le|\ell_\xi(z)|\le\frac54L_\xi.
\]
The higher logarithmic derivatives satisfy the weighted bounds needed
for \eqref{eq:H-relative}: indeed,
$k\log k\le CL_\xi^2$ and $k^{j/2}\le L_\xi^j$ for $j\ge3$.
Moreover,
\[
 |\nabla\arg G_{\xi,\eps_\xi}|
 \le\frac2{\eps_\xi}+C=o(L_\xi).
\]
The reverse triangle inequality and \eqref{eq:holomorphic-phase}
therefore give
\begin{equation}\label{eq:phase-A}
 |\nabla\arg F|
 \ge|\ell_\xi|-|\nabla\arg G_{\xi,\eps_\xi}|
 \ge\frac12L_\xi.
\end{equation}
Here the holomorphic phase dominates. No alignment of the two
phase gradients is required; even opposite directions are allowed.

\smallskip\noindent\emph{Case B: $a_\xi<k^{2/3}$.}
We now shrink the radius so that the replacement phase dominates.
Set $\eps_\xi=\eps_1k^{-2/3}$ and
$L_\xi=\eps_\xi^{-1}$. On the smaller disk,
\[
 |\ell_\xi(z)|
 \le k^{2/3}+C\eps_1k^{1/3}\log k\le2k^{2/3}
\]
for all sufficiently large $k$. The same weighted derivative
bounds again imply \eqref{eq:H-relative} with $L=L_\xi$.
On the enlarged ring, \eqref{eq:Gphase} gives
\begin{align*}
 |\nabla\arg F|
 &\ge\frac1{3\eps_\xi}-C-2k^{2/3}\\
 &=\left(\frac13-2\eps_1\right)L_\xi-C.
\end{align*}
Choose $\eps_1<1/24$, independently of $k$ and $\xi$, and then
increase $k_0$ so that $C\le L_\xi/8$. We obtain
\begin{equation}\label{eq:phase-B}
 |\nabla\arg F|\ge L_\xi/8.
\end{equation}

\smallskip\noindent\emph{Residual estimates.}
Since $\eps_\xi^{-1}\le CL_\xi$, the relative derivative
bounds \eqref{eq:Gder} and \eqref{eq:H-relative} give
\eqref{eq:patchF}. To estimate the residual, the three relative
costs in \eqref{eq:prod} are
\[
 \frac{|\pb^2G|}{|G|}\le C\eps_\xi^{-2},\qquad
 \frac{|\pa\pb^2G|}{|G|}\le C\eps_\xi^{-3},\qquad
 \frac{|\pa^2\pb^2G|}{|G|}\le C\eps_\xi^{-4}.
\]
After $j\le5$ additional real derivatives, the product rule yields
\begin{equation}\label{eq:residual-budget}
 \begin{split}
 |D^j\Delta^2(H_\xi G)|
 \le C_j L_\xi^j\bigl(L_\xi^2\eps_\xi^{-2}
                +L_\xi\eps_\xi^{-3}+\eps_\xi^{-4}\bigr)|H_\xi G|.
 \end{split}
\end{equation}
Each extra derivative costs at most $CL_\xi$; the largest derivative
orders needed are seven for $H_\xi$ and nine for $G$.

In Case A, $L_\xi\eps_\xi\to\infty$, so the first term
in \eqref{eq:residual-budget} dominates. We take
\begin{equation}\label{eq:budget-A}
 D_\xi=L_\xi^2\eps_\xi^{-2}.
\end{equation}
Since $L_\xi\le Ck$,
\[
 D_\xi\le Ck^3,\qquad
 D_\xi/L_\xi^4\le Ck/L_\xi^2\le Ck^{-1/3}.
\]
In Case B all three terms have order $L_\xi^4$, so take
\begin{equation}\label{eq:budget-B}
 D_\xi=L_\xi^4=\eps_1^{-4}k^{8/3}\le Ck^3.
\end{equation}
Since $\eps_1$ is fixed, we in fact have $D_\xi/k^3\to0$
in this case. Thus the smaller radius still leaves room within
the potential budget. Both choices give \eqref{eq:patchE} and
the remaining bounds in \eqref{eq:patchphase}.

Finally, exact biharmonicity confines the residual to the stated
closed transition ring. Its distance from $\pa\Omega_\xi$ is at least
$\eps_\xi/2$. In Case A,
$L_\xi\eps_\xi\ge\eps_0k^{1/6}$; in Case B,
$L_\xi\eps_\xi=1$. A fixed $b>0$ therefore gives the required
margin for all sufficiently large $k$.
\end{proof}

\subsection{The real annular connection}
\label{sec:annular}

All regions carrying a residual now satisfy the hypotheses for
real-part correction, except the inner collar, where the real part
already stays away from zero. Because these regions are disjoint
and the corrections preserve their boundary collars, the local
potentials assemble without any loss proportional to the number
of roots.

\begin{proposition}\label{prop:annular}
For every sufficiently large square integer $k=q^2$, there exist
$f_k,Q_k\in C^\infty(\R^2;\R)$ such that
\begin{equation}\label{eq:annular}
 \Delta^2f_k+Q_kf_k=0,\qquad
 Q_k\in C_c^\infty(\{1<|z|<4\};\R),\qquad
 \norm{Q_k}_\infty\le Ck^3,
\end{equation}
and
\begin{equation}\label{eq:annmodels}
 \begin{aligned}
 f_k(z)&=A_k|z|^2&&(|z|\le1),\\
 f_k(z)&=\Rea(|z|^2z^{-k})&&(|z|\ge4),
 \end{aligned}
 \qquad A_k=P(0)^{-1}\in\R\setminus\{0\}.
\end{equation}
\end{proposition}

\begin{proof}
Start with $\Rea F$, where $F$ is the global smooth function from
Section~\ref{sec:adaptive}. On each $\Omega_\xi$, apply
Lemma~\ref{lem:repair} using Proposition~\ref{prop:patch}.
The estimates for $\Delta^2F$ imply those for
$E=\Delta^2\Rea F$, and $\supp E$ is contained in the same
transition ring. All constants in the correction lemma are uniform
in $\xi,k$.

On the outer enlarged collar $\Omega_+$, apply the same lemma with
\begin{equation}\label{eq:outer-LD}
 L=k,\qquad D=k^2q\vartheta^q.
\end{equation}
Equations~\eqref{eq:collarF}--\eqref{eq:outerphase} give its
hypotheses. In particular, $D/L^4=q\vartheta^q/k^2$ is bounded,
and the residual support $\mathcal A_+$ has a fixed positive
boundary margin. The lower bound on $\nabla\arg F^0$ holds on the
whole enlarged collar.

The root disks are pairwise disjoint and avoid both
$\Omega_-$ and $\Omega_+$. Each correction is supported in
the original residual support, compactly inside its enlarged
region. Hence the corrected functions agree with $\Rea F$
on open neighborhoods of every gluing boundary and define a
single $f_k\in C^\infty(\R^2;\R)$. Both exact models in
\eqref{eq:annmodels} are preserved.

On $\Omega_-$, no correction is made. By \eqref{eq:inner-real},
we may define
\[
 Q_k=-\frac{\Delta^2\Rea F^0}{\Rea F^0},\qquad
 |Q_k|\le Ck^2q\vartheta^q.
\]
This is smooth and has support in $\mathcal A_-$.
On every corrected region, use the smooth real potential supplied
by Lemma~\ref{lem:repair}. On the remaining open biharmonic
regions, set $Q_k=0$.

The potentials agree on overlaps because each vanishes in an
open neighborhood of the boundary of its enlarged region. In
particular, the core $|z-\xi|<\eps_\xi$ retains the exact
formula
\[
 f_k(z)=\Rea\left(
     \frac{\bar\xi}{\eps_\xi^2}H_\xi(z)(\bar z-\bar\xi)\right),
 \qquad \Delta^2f_k=0.
\]
Setting $Q_k=0$ in this open disk verifies the equation even
at a critical zero of the real part; no division is performed
there. The original replacements also agree on open overlap
regions, so differentiation introduces no interface distributions.

The potential support is contained in
\begin{equation}\label{eq:annular-support}
 \mathcal A_-\cup\mathcal A_+
 \cup\bigcup_{\xi\in\Lambda}
      \{\eps_\xi\le|z-\xi|\le2\eps_\xi\},
\end{equation}
a compact subset of $1<|z|<4$. At each point at most one of these
local potentials is nonzero. Hence
\[
 \norm{Q_k}_\infty
 \le\max\left\{
       \sup_{\xi\in\Lambda}CD_\xi,\ Ck^2q\vartheta^q\right\}
 \le Ck^3.
\]
The estimate is independent of the number of roots. This proves
\eqref{eq:annular}.
\end{proof}


\section{Kelvin inversion and normalization}
\label{sec:kelvin}

The annular construction has supplied a real pair with an exact
model at each end and a potential supported between them. Kelvin
inversion exchanges these models, producing a high-order zero
near the origin and a nonzero constant near infinity. Translation,
normalization at a point of maximum modulus, and dilation will then
place a unit observation ball at exactly the prescribed distance.

\subsection{The transformed equation and its exact models}

For $y\ne0$, let $I(y)=y/|y|^2$. The planar biharmonic Kelvin
identity is
\begin{equation}\label{eq:Kelvin}
 \Delta_y^2\bigl(|y|^2f(I(y))\bigr)
 =|y|^{-6}(\Delta^2f)(I(y)).
\end{equation}
For completeness, in polar coordinates put $D_r=r\pa_r$. Then
\[
 \Delta^2=r^{-4}
     \bigl((D_r-2)^2+\pa_\theta^2\bigr)
     \bigl(D_r^2+\pa_\theta^2\bigr).
\]
With $s=r^{-1}$, the operator $D_r$ acting on
$r^2f(s,\theta)$ corresponds to $2-D_s$ acting on $f$.
This substitution exchanges the two commuting factors. Hence
\[
 \Delta_y^2(r^2f(s,\theta))
 =r^{-2}(D_s^2+\pa_\theta^2)
              ((D_s-2)^2+\pa_\theta^2)f
 =r^{-2}s^4\Delta_z^2f,
\]
which proves \eqref{eq:Kelvin}. The calculation is local away
from the origin and requires only $f\in C^4$.

Apply it to Proposition~\ref{prop:annular}, and define
\begin{equation}\label{eq:kelvin-pair}
 g_k(y)=|y|^2f_k(I(y)),\qquad
 \widehat Q_k(y)=|y|^{-8}Q_k(I(y))\qquad(y\ne0).
\end{equation}
The exact models give
\begin{equation}\label{eq:gmodels}
 g_k(y)=\Rea(\bar y^{\,k})\quad(|y|\le1/4),\qquad
 g_k(y)=A_k\quad(|y|\ge1).
\end{equation}
Define $\widehat Q_k=0$ at zero. Since it vanishes in a whole
neighborhood of zero, both functions extend smoothly there, and
\begin{equation}\label{eq:kelvin-budget}
 \Delta^2g_k+\widehat Q_kg_k=0,\qquad
 \supp\widehat Q_k\Subset\{1/4<|y|<1\},\qquad
 \norm{\widehat Q_k}_\infty\le C_Kk^3.
\end{equation}
On the transformed support, $|y|^{-8}\le4^8$, so we may
take $C_K=4^8C$, with $C$ as in \eqref{eq:annular}.
The support restriction is essential: Kelvin inversion does not
preserve boundedness for an arbitrary bounded potential. The exact
polynomial near zero also rules out any point-supported
distribution at the origin.

Let $M_k=\norm{g_k}_\infty$. Evaluation at $y=(1/4,0)$ yields
$M_k\ge4^{-k}$. There is a point $y_k$ such that
\begin{equation}\label{eq:max}
 1/4\le|y_k|\le1,\qquad |g_k(y_k)|=M_k.
\end{equation}
Indeed, $|g_k(y)|\le|y|^k<4^{-k}$ for $|y|<1/4$, while
for $|y|\ge1$ the function is constant and assumes that same
value on $|y|=1$. The supremum is therefore attained on the
compact annulus in \eqref{eq:max}. This uses the exact models
and compactness, not a maximum principle.

\subsection{Dilation and completion of the proof}

\begin{proof}[Proof of Theorem~\ref{thm:main}]
We first work in dimension two. Fix
$L_0\ge\max\{1,C_K^{1/4}\}$. Given a sufficiently large
real number $R$, choose the degree and dilation by
\begin{equation}\label{eq:parameters}
 q=\left\lfloor(R/L_0)^{2/3}\right\rfloor,\qquad
 k=q^2,\qquad \lambda=R/|y_k|.
\end{equation}
For $(R/L_0)^{2/3}\ge2$,
\begin{equation}\label{eq:k-comparison}
 \frac{1}{4L_0^{4/3}}R^{4/3}\le k
       \le\frac{1}{L_0^{4/3}}R^{4/3}.
\end{equation}
Also $R\le\lambda\le4R$ and
\begin{equation}\label{eq:dilation-budget}
 L_0k^{3/4}=L_0q^{3/2}\le R\le\lambda.
\end{equation}
Define
\begin{equation}\label{eq:final-pair}
 \begin{aligned}
 u_R(x)&=\frac{g_k(y_k+x/\lambda)}{g_k(y_k)},\\
 V_R(x)&=\lambda^{-4}\widehat Q_k(y_k+x/\lambda),
 \qquad x_R=-\lambda y_k.
 \end{aligned}
\end{equation}
The denominator is nonzero; dividing by its actual value fixes
the sign at the origin. By \eqref{eq:kelvin-budget} and
\eqref{eq:dilation-budget},
\[
 \Delta^2u_R+V_Ru_R=0,\qquad
 \norm{V_R}_\infty\le C_Kk^3\lambda^{-4}\le1.
\]
The choice of $y_k$ gives $u_R(0)=\norm{u_R}_\infty=1$, and
$|x_R|=R$. The planar potential is smooth and compactly supported.

The chosen translation sends the observation ball into the
polynomial region. More precisely, for $x=x_R+v$ with $|v|<1$,
the argument of $g_k$ is $v/\lambda$; when $\lambda>4$,
\eqref{eq:gmodels} applies and gives
\begin{equation}\label{eq:ball-estimate}
 |u_R(x)|\le\frac{\lambda^{-k}}{M_k}
 \le(4/\lambda)^k\le(4/R)^k.
\end{equation}
If $R\ge16$, then $\log(R/4)\ge\tfrac12\log R$.
Together with \eqref{eq:k-comparison}, this proves
\eqref{eq:example-smallness}, for example with
$c=(8L_0^{4/3})^{-1}$.

For $n>2$, set
\[
 u_R^{(n)}(x)=u_R^{(2)}(x_1,x_2),\qquad
 V_R^{(n)}(x)=V_R^{(2)}(x_1,x_2),
\]
and append zero coordinates to $x_R$. Then
$\Delta_n^2u_R^{(n)}=\Delta_2^2u_R^{(2)}$, and the global norms
are unchanged. The projection of $B_1^{(n)}((x_R,0))$ onto the first two
coordinates is exactly $B_1^{(2)}(x_R)$. Continuity and independence
of the remaining coordinates therefore give equality of the
unit-ball supremum norms, proving the asserted smallness in every
dimension. The higher-dimensional support is generally cylindrical,
not compact.

For each fixed $R$, the planar solution is a nonzero constant
outside a sufficiently large ball, by the exterior model in
\eqref{eq:gmodels}. This completes the proof.
\end{proof}


\section{Concluding remarks}
\label{sec:scope}

\begin{remark}[Comparison with the second-order equation]\label{rem:second-order}
The fourth-order mechanism is not merely the availability of a
local jet correction. For holomorphic $H$, the identity
\[
 \Delta^2\bigl(H(z)(\bar z-\bar\xi)\bigr)=0
\]
keeps the root cores exactly biharmonic. By contrast, the
second-order identity is
\[
 \Delta\bigl(H(z)(\bar z-\bar\xi)\bigr)=4H'(z),
\]
whose right-hand side need not vanish, even on the zero set of
the real part. The same root-core construction therefore does not
yield a bounded real potential for the second-order equation.
The local identity $\Delta(h^2a)|_{h=0}=2a|\nabla h|^2$
does provide a second-order cancellation, but it does not make
these cores harmonic. What matters here is that possible critical
zeros lie in open regions where the biharmonic residual vanishes
identically. This explains the present construction; it does not
exclude other second-order methods.
\end{remark}

\begin{remark}[Finite-order zeros and the logarithm]
Near zero, the transformed solution is exactly
$g_k(y)=\Rea(\bar y^{\,k})$. Its leading jet has degree $k$,
so its vanishing order is finite and exactly $k$; the construction
produces neither an infinite-order zero nor a nontrivial compactly
supported solution. For the final solution,
\[
 u_R(x_R+v)
 =\frac{\lambda^{-k}\Rea(\bar v^{\,k})}{g_k(y_k)}
 \qquad(|v|<\lambda/4)
\]
in dimension two. The logarithm is thus a consequence of dilation at a zero of
high finite order: $\lambda^{-k}=\exp(-k\log\lambda)$,
with $k\asymp R^{4/3}$ and $\lambda\asymp R$. The
normalization uses attainment of a maximum on a compact annulus,
not a maximum principle for the biharmonic operator.
\end{remark}

\begin{remark}[Order of the parameter choices]\label{rem:constants}
The fixed radii, cutoffs and separation constant determine $\eps_0$;
then $\eps_1$ is fixed and the uniform large-degree threshold is
chosen. The resulting uniform constants in
Lemma~\ref{lem:repair} determine $\delta$, after taking the minimum
of the root and outer-collar margin constants and the maximum of
their derivative constants. The correction estimates determine
$C_K$, which determines $L_0$. Only then is $R_*$ fixed to
ensure that \eqref{eq:parameters} gives an admissible degree and that
the final ball and logarithm estimates apply. In particular, neither
$\delta$ nor the fixed dilation constant depends on $k$ or $R$.
\end{remark}

\begin{remark}[Scalar classes and scope]
The potentials in the final examples are scalar and real.
Complex functions enter only as auxiliary devices in constructing
that real pair. The inclusions
\[
 \mathcal C_n^{\R,\R}(C_0)
 \subset\mathcal C_n^{\C,\R}(C_0)
 \subset\mathcal C_n^{\C,\C}(C_0)
\]
give the complex-class conclusion of Corollary~\ref{cor:sharpness}
without a separate construction. The result concerns
$\Delta^2+V$ and does not extend, as stated, to arbitrary
fourth-order elliptic principal parts, fixed-sign potentials, or
classes with additional uniform regularity or support restrictions.
Nor does it determine the optimal all-direction decay of a single
fixed real solution.
\end{remark}

\begin{thebibliography}{99}

\bibitem{Bakri12}
L.~Bakri,
\emph{Quantitative uniqueness for Schr\"odinger operator},
Indiana Univ. Math. J. \textbf{61} (2012), no.~4, 1565--1580.

\bibitem{Bakri13}
L.~Bakri,
\emph{Carleman estimates for the Schr\"odinger operator.
Applications to quantitative uniqueness},
Comm. Partial Differential Equations \textbf{38} (2013), no.~1, 69--91.

\bibitem{BK}
J.~Bourgain and C.~E.~Kenig,
\emph{On localization in the continuous Anderson--Bernoulli model
in higher dimension},
Invent. Math. \textbf{161} (2005), no.~2, 389--426.

\bibitem{CG}
F.~Colombini and C.~Grammatico,
\emph{Some remarks on strong unique continuation for the Laplace
operator and its powers},
Comm. Partial Differential Equations \textbf{24} (1999), nos.~5--6,
1079--1094.

\bibitem{CS}
J.~Cruz-Sampedro,
\emph{Unique continuation at infinity of solutions to Schr\"odinger
equations with complex-valued potentials},
Proc. Edinburgh Math. Soc. (2) \textbf{42} (1999), no.~1, 143--153.

\bibitem{Davey}
B.~Davey,
\emph{Some quantitative unique continuation results for
eigenfunctions of the magnetic Schr\"odinger operator},
Comm. Partial Differential Equations \textbf{39} (2014), no.~5,
876--945.

\bibitem{DKW}
B.~Davey, C.~Kenig, and J.-N.~Wang,
\emph{The Landis conjecture for variable coefficient second-order
elliptic PDEs},
Trans. Amer. Math. Soc. \textbf{369} (2017), no.~11, 8209--8237.

\bibitem{DF}
H.~Donnelly and C.~Fefferman,
\emph{Nodal sets of eigenfunctions on Riemannian manifolds},
Invent. Math. \textbf{93} (1988), no.~1, 161--183.

\bibitem{FGJL}
S.~Feng, Z.~Gan, R.~Jiang, and F.-H.~Lin,
\emph{Vanishing orders of Dirichlet solutions to the Schr\"odinger
equation in dimensions three and higher},
preprint (2026),
\href{https://arxiv.org/abs/2610.01615v1}{arXiv:2610.01615v1}.

\bibitem{FI}
R.~L.~Frank and P.~Ivanisvili,
\emph{Counterexamples to the Landis conjecture in dimensions three
and higher},
preprint (2026),
\href{https://arxiv.org/abs/2608.00802v1}{arXiv:2608.00802v1}.

\bibitem{FZZ}
X.~Fu, X.~Zhang, and E.~Zuazua,
\emph{On the optimality of some observability inequalities for plate
systems with potentials}, in \emph{Phase Space Analysis of Partial
Differential Equations} (A.~Bove, F.~Colombini, and D.~Del Santo, eds.),
Progr. Nonlinear Differential Equations Appl., vol.~69,
Birkh\"auser, Boston, MA, 2006, pp.~117--132.

\bibitem{GL86}
N.~Garofalo and F.-H.~Lin,
\emph{Monotonicity properties of variational integrals, $A_p$ weights
and unique continuation},
Indiana Univ. Math. J. \textbf{35} (1986), no.~2, 245--268.

\bibitem{GL87}
N.~Garofalo and F.-H.~Lin,
\emph{Unique continuation for elliptic operators: a
geometric-variational approach},
Comm. Pure Appl. Math. \textbf{40} (1987), no.~3, 347--366.

\bibitem{HWZ}
S.~Huang, M.~Wang, and Q.~Zheng,
\emph{Quantitative uniqueness of some higher order elliptic equations},
J. Math. Anal. Appl. \textbf{444} (2016), no.~1, 326--339.
Author version:
\href{https://arxiv.org/abs/1505.05228v1}{arXiv:1505.05228v1}.

\bibitem{JK}
D.~Jerison and C.~E.~Kenig,
\emph{Unique continuation and absence of positive eigenvalues for
Schr\"odinger operators},
Ann. of Math. (2) \textbf{121} (1985), no.~3, 463--488.

\bibitem{Kenig}
C.~E.~Kenig,
\emph{Some recent applications of unique continuation}, in
\emph{Recent Developments in Nonlinear Partial Differential Equations},
Contemp. Math., vol.~439, Amer. Math. Soc., Providence, RI, 2007,
pp.~25--56.

\bibitem{KSW}
C.~Kenig, L.~Silvestre, and J.-N.~Wang,
\emph{On Landis' conjecture in the plane},
Comm. Partial Differential Equations \textbf{40} (2015), no.~4,
766--789.

\bibitem{Kukavica}
I.~Kukavica,
\emph{Quantitative uniqueness for second-order elliptic operators},
Duke Math. J. \textbf{91} (1998), no.~2, 225--240.

\bibitem{Lin07}
C.-L.~Lin,
\emph{Strong unique continuation for $m$-th powers of a Laplacian
operator with singular coefficients},
Proc. Amer. Math. Soc. \textbf{135} (2007), no.~2, 569--578.

\bibitem{LNW}
C.-L.~Lin, S.~Nagayasu, and J.-N.~Wang,
\emph{Quantitative uniqueness for the power of the Laplacian with
singular coefficients},
Ann. Sc. Norm. Super. Pisa Cl. Sci. (5) \textbf{10} (2011), no.~3,
513--529.

\bibitem{LW}
C.-L.~Lin and J.-N.~Wang,
\emph{Quantitative uniqueness estimates for the general second order
elliptic equations},
J. Funct. Anal. \textbf{266} (2014), no.~8, 5108--5125.

\bibitem{LTY}
H.~Liu, L.~Tian, and X.~Yang,
\emph{Quantitative unique continuation property for solutions to a
bi-Laplacian equation with a potential},
J. Lond. Math. Soc. (2) \textbf{112} (2025), no.~2, e70265.

\bibitem{LMNN}
A.~Logunov, E.~Malinnikova, N.~Nadirashvili, and F.~Nazarov,
\emph{The Landis conjecture on exponential decay},
Invent. Math. \textbf{241} (2025), 465--508.

\bibitem{Meshkov}
V.~Z.~Meshkov,
\emph{On the possible rate of decay at infinity of solutions of
second order partial differential equations},
Math. USSR-Sb. \textbf{72} (1992), no.~2, 343--361.

\bibitem{Zhu16}
J.~Zhu,
\emph{Quantitative uniqueness of elliptic equations},
Amer. J. Math. \textbf{138} (2016), no.~3, 733--762.

\bibitem{Zhu18}
J.~Zhu,
\emph{Quantitative unique continuation of solutions to higher order
elliptic equations with singular coefficients},
Calc. Var. Partial Differential Equations \textbf{57} (2018), no.~2,
Paper No.~58, 35~pp.

\end{thebibliography}
\end{document}
